\section{Related Work}

\subsection{Text-to-Image Evaluation}
Most current image generation metrics were not designed for reference-conditioned multi-subject identity preservation. Global embedding metrics such as CLIP \cite{radford2021learning} and DINOv2 \cite{oquab2023dinov2} are effective when the dominant question is whether an image matches a caption at a coarse semantic level, but they tend to average away localized identity failures. General reward and preference models, such as ImageReward \cite{xu2023imagereward} and PickScore \cite{kirstain2023pick}, can provide useful aesthetic or holistic judgments, yet they usually remain black-box signals and often prioritize visual appeal over strict identity fidelity. Our work focuses on the regime where these shortcomings become critical: multiple reference identities, non-trivial interactions, and identity-preserving generation under the same prompt.

\subsection{Compositional and VQA-Style Benchmarks}
Prior composition-oriented benchmarks like T2I-CompBench \cite{huang2023t2i} demonstrate that text-to-image models struggle when multiple attributes or spatial relations must be satisfied at once. However, most of these settings assume text-only specification rather than explicit reference subjects. This distinction matters. In multi-subject personalization, the task is not only to realize a caption, but to correctly bind the right visual identity to the right local attributes and interactions. A benchmark that ignores reference images cannot fully evaluate identity collapse, clone failure, or cross-subject appearance leakage.

\subsection{Subject-Driven and Personalized Generation}
Subject-driven generation methods are typically evaluated with single-subject identity preservation metrics such as ArcFace \cite{deng2019arcface} or limited qualitative examples. The multi-subject regime introduces a different failure geometry, as noted in our prior work on multi-subject evaluation \cite{chen2024multisubject}. Identity preservation must scale with subject count, and the evaluator must remain discriminative as the scene grows denser. This gap motivates a benchmark that directly stress-tests scalability rather than treating multi-subject generation as a minor extension of the single-subject case.

\subsection{Learned Evaluators and Diagnostic Metrics}
Learned metrics are especially promising when existing handcrafted or global embedding metrics fail to reflect human judgment. Our work follows this spirit but differs in two important ways. First, AnchorScore is explicitly reference-conditioned and pairwise trained under shared prompts. Second, it produces both a scalar preference signal and interpretable diagnostic dimensions. This makes the metric useful not only for ranking models, but also for diagnosing why they fail.
